\section*{Conclusion}
\label{conclusion}

This study asked how much of a voice can be reconstructed from material that was
never recorded for that purpose: $42$ WhatsApp voice messages, half an hour in
total, from a patient who has since lost her speech. Three commercial zero-shot
services were compared against her own recordings as a fourth condition, using
objective metrics, a naturalness test and a voice Turing test with $30$ raters.
Two of the three systems are rated as natural as the genuine recordings, none
passes as the patient in a direct comparison, and S2 Pro is recommended for
deployment (Section~\ref{model_selection}). Legacy material is therefore
sufficient for an assistive voice, if not for a replica of the original.

Three lessons carry beyond this case. First, zero-shot cloning from a reference
of a few seconds matched the fine-tuned model in a listening comparison at a
fraction of the compute, which made the fine-tuning path unnecessary. Second,
preprocessing was not made redundant by that change, as was expected at the
time. DeepFilterNet removed the background noise from the phone recordings
without audible damage, and the cleaned reference is what all three services
were enrolled with; yet the same denoising produced the digitally silent pauses
that the listening test then had to neutralise. A preprocessing step can solve
one problem and create another further down the pipeline. Third, the objective
metrics pointed at the wrong system, and only the listening test made that
visible. That the same ease of cloning has a second reading is set out in the
Discussion.

The next step is deployment. The pipeline built here becomes the backend of an
application in which a user uploads a recording, the file is cleaned in the
background and a clone is created from it. Two points remain open before that.
Self-hosting, one of the stated requirements, was not tested, since all three
systems were used through their providers' APIs. And a follow-up study would
need several speakers, a hearing screening, and raters whose age is not tied to
their familiarity with the speaker.
